\documentclass[11pt,letterpaper]{article}
\usepackage[T1]{fontenc}
\usepackage{lmodern,amsmath,amssymb,amsthm,booktabs}
\usepackage[margin=1.0in,top=.84in,bottom=.87in,headheight=14pt,headsep=15pt]{geometry}
\usepackage{xcolor,fancyhdr,microtype}
\usepackage[colorlinks=true,linkcolor=accent,urlcolor=accent,citecolor=accent,pdfauthor={Research note},pdftitle={When do disjoint products of small primes cover every residue tuple?}]{hyperref}
\definecolor{accent}{RGB}{17,84,95}
\newtheorem{theorem}{Theorem}
\newtheorem{lemma}[theorem]{Lemma}
\newtheorem{corollary}[theorem]{Corollary}
\newcommand{\G}{G}
\newcommand{\Pp}{\mathcal P}
\newcommand{\E}{\mathbb E}
\newcommand{\F}{\mathbb F}
\newcommand{\pr}{\mathbb P}
\newcommand{\abs}[1]{\lvert#1\rvert}
\newcommand{\bias}{\operatorname{bias}}
\newcommand{\one}{\mathbf 1}
\newcommand{\srcbase}{https://github.com/VladimirReshetnikov/ProveIt/blob/b6253012cf30c80a989c871906e097a27f23a27c/Analysis/GeneralizedRiemannHypothesis}
\newcommand{\repo}[2]{\href{\srcbase/#1}{#2}}
\setlength{\parindent}{1.15em}
\setlength{\parskip}{2pt plus .3pt}
\setlength{\emergencystretch}{2em}
\pagestyle{fancy}\fancyhf{}
\fancyhead[L]{\footnotesize\sffamily\color{accent} Disjoint products of small primes}
\fancyhead[R]{\footnotesize\sffamily Arithmetic combinatorics}
\fancyfoot[C]{\small\thepage}
\renewcommand{\headrulewidth}{.3pt}
\begin{document}
\thispagestyle{plain}
{\noindent\sffamily\small\bfseries\color{accent} RESEARCH NOTE \hfill 27 NOVEMBER 2026\par}
\vspace{12pt}
{\raggedright\noindent\Huge\bfseries When do disjoint products of small primes cover every residue tuple?\par}
\vspace{8pt}
{\noindent\large\color{accent} Exact quadratic and cubic obstructions after arbitrary deletions\par}
\vspace{13pt}
\noindent\textbf{Abstract.}
Fix \(r\) target residues modulo \(q\). From the primes up to \(C(\log q)^2\) retain any fixed fraction \(a>1/4\). Form \(r\) pairwise coprime squarefree products with prescribed numbers of prime factors, leaving a prescribed reserve. When every quota is at least \(B\log q\), we give necessary and sufficient conditions for any target tuple to occur.
For a quadratic character, all tuples occur exactly as far as that character is concerned when at least \(r\) retained primes have its less frequent sign. For a cubic character, let \(u,v\) count the two classes outside a most frequent class. Its exact full-coverage threshold is \(u+v+\min(u,v)\ge 2r\). These tests, over all characters of orders two and three, suffice together. We also compute the transition to pointwise equidistribution and show that the density cutoff \(1/4\) for this coverage classification is sharp. At every fixed positive retention density, a single list of \(O(R\log Q)\) primes also gives relative error \(Q^{-A}\) within each occupied quotient fibre for up to \(R\) independently chosen moduli at most \(Q\). The range \(R\asymp\log Q/\log\log Q\) is optimal in order for these quotas and this prime cutoff. The arithmetic input is Dirichlet GRH as stated in the vendored Lean development. The allocation and counting arguments are unconditional.

\section{The question}
How can a few extra primes change which products are possible? Suppose that every retained prime is a quadratic residue modulo an odd prime, except for one nonresidue. One can put that prime in either of two products, or leave it unused. Each product separately may attain either sign. They can never both have negative sign. A second nonresidue removes this joint obstruction. Even then, the four sign patterns generally have different probabilities.

The source article asks whether the distance from a quadratic fibre can replace a fixed density margin above one half in its theorem on squarefree products \cite{article}. A companion report in the same repository gives a finite quotient governing several disjoint products and a dynamic program for their counts \cite{factor}. Our contribution is an explicit description of \emph{each target} above quarter density, a sharp coverage threshold, and an if-and-only-if transition to uniform counts. The expanded article also asks how many factors its relative-count theorem permits below majority retention \cite{merged}. Theorems~\ref{th:growing} and~\ref{th:mixed} give the maximal order, including when each factor has its own modulus. A repair of the integer table of block sizes and character values keeps the Fourier error smaller than every positive projected mass as the number of blocks grows.

\paragraph{Notation.}
Fix an integer \(r\ge1\), a real \(1/4<a\le1\), and an integer modulus \(q\ge3\).
Write \(G=(\mathbb Z/q\mathbb Z)^\times\), \(L=\log q\), \(y=CL^2\), and
\[
 \Pp=\{p\le y:p\text{ prime},\ p\nmid q\},\quad N=\abs{\Pp},
 \quad S\subseteq\Pp,\quad m=\abs S\ge aN .
\]
Choose a uniform ordered partition \(S=S_0\sqcup\cdots\sqcup S_r\), with prescribed integer sizes \(n_i\), \(\sum_{i=0}^r n_i=m\), and write
\[
 X=(X_1,\ldots,X_r),\qquad X_j=\prod_{p\in S_j}p\bmod q,\qquad
 h=\min_{0\le i\le r} n_i,\qquad t=h/m.
\]
Block zero is the reserve. There are \(m!/\prod n_i!\) equiprobable partitions. Unique factorization identifies each partition with its ordered tuple of actual squarefree integers. The constants \(B,C\) below depend only on \(a,r\). All arithmetic conclusions apply to every sufficiently large \(q\) and every admissible \(S,(n_i)\) with \(h\ge B L\). A character means a character of \(G\).

\section{Closed tests for coverage and uniformity}
For a character \(\chi\) of order \(d=2\) or \(3\), choose a most frequent value \(b\) on \(S\), and write \(s_\chi=\#\{p\in S:\chi(p)\ne b\}\). For a specified target \(g=(g_1,\dots,g_r)\), its corrected character values are \(w_j=\chi(g_j)b^{-n_j}\). Fix \(\omega=e^{2\pi i/3}\).

\begin{theorem}[Exact support]\label{th:support}
Using the Dirichlet GRH input, a target \(g\in G^r\) occurs if and only if it passes both of the following tests for every character of the indicated order.
\begin{itemize}
\item For order two, the number of \(j\) with \(w_j=-1\) is at most \(s_\chi\).
\item For order three, put \(u=\#\{p:\chi(p)=b\omega,\ p\in S\}\), \(v=\#\{p:\chi(p)=b\omega^{-1},\ p\in S\}\), and let \(i,j\) count respectively the corrected entries \(\omega,\omega^{-1}\). Then
\begin{equation}\label{eq:polygon}
             2u+v\ \ge\ 2i+j,\qquad u+2v\ \ge\ i+2j .
\end{equation}
\end{itemize}
In particular, \(X\) covers all of \(G^r\) if and only if
\begin{equation}\label{eq:coverage}
  s_\chi\ge r\quad(\operatorname{ord}\chi=2),\qquad
  u+v+\min(u,v)\ge 2r\quad(\operatorname{ord}\chi=3).
\end{equation}
For \(a>1/3\) the quadratic tests alone suffice. Any tie for \(b\) can be resolved arbitrarily.
\end{theorem}

For a pointwise measure of uniformity, put
\[
 D_q=\max_{g\in G^r}\left|\abs{G}^{\,r}\pr(X=g)-1\right|,\quad
 B_2(t)=1-2t,\qquad B_3(t)=\sqrt{1-3t(1-t)}.
\]
Let \(\mathcal M_q=\max_\chi B_{\operatorname{ord}\chi}(t)^{s_\chi}\), over all order-two and order-three characters, interpreting \(0^0=1\) and an empty maximum as zero.

\begin{theorem}[Pointwise mixing]\label{th:mix}
Under the same assumptions, uniformly in the retained primes and the quotas,
\[
 \mathcal M_q-O_{a,r}\!\left(\frac{\log(2h)}h\right)
 \ \le D_q\ \le\
 (3^r-1)\mathcal M_q+O_{a,r}\!\left(\frac{\log(2h)}h\right).
\]
In particular \(D_q\to0\) along an admissible sequence if and only if \(\mathcal M_q\to0\). For \(a>1/3\) one may use just the quadratic characters and replace \(3^r-1\) by \(2^r-1\). For \(r=1,\ a>1/3\) this gives the sharper identity
\[
 D_q=(1-2h/m)^{\,\min_{\operatorname{ord}\chi=2}s_\chi}
             +O_a\!\left(\frac{\log(2h)}h\right).
\]
\end{theorem}

When \(r\ge2\), \(t\le 1/3\), so uniformity is equivalent to \(t\min_\chi s_\chi\to\infty\) over the relevant characters. Indeed \(-\log B_d(t)\) is bounded above and below by positive multiples of \(t\) on this range. For one product the quadratic exception at \(t=1/2\) is real: exactly half of the partitions place a single exceptional prime in that product. Its sign is then perfectly balanced. A cubic exception has \(B_3(1/2)=1/2\), so boundedly many such primes can retain a detectable bias.

\section{From Dirichlet GRH to a quotient of order at most three}
\paragraph{The arithmetic input.}
We use only the following prime-pool estimate. Given \(\epsilon>0\), one can take \(C=C(\epsilon)\) so that, for every sufficiently large \(q\),
\begin{equation}\label{eq:pool}
  N\sim y/\log y,\qquad
  \max_{\psi\ne1}\left|\frac 1N\sum_{p\in\Pp}\psi(p)\right|\le\epsilon .
\end{equation}
This is \texttt{dir:lem:bias} of \cite{article}. At the endpoint \(y=C(\log q)^2\), its proof uses the once-integrated GRH explicit-formula estimate
\[
 \sum_{n\le x}(x-n)\Lambda(n)\psi(n)\ll x^{3/2}\log(2q)
 \quad(\psi\ne1)
\]
as in \cite[Ch.~13]{MV}. It applies also to imprimitive characters by accounting for the omitted Euler factors. For clarity, subtract this estimate at \((1-\eta)y\) from that at \(y\), and divide the \emph{difference of sums} by \(\eta y\). The resulting weight is one up to \((1-\eta)y\) and decreases linearly to zero. Restoring the endpoint costs at most \((\eta+o(1))y\) by the prime number theorem. Higher prime powers cost \(O(\sqrt y\log^2 y)\), and replacing \(\log p\) by \(\log y\) costs \(o(y)\) on summation. There are at most \(L/\log 2\) excluded prime divisors. Division by \(N\log y\sim y\) yields \(\eta+O(1/(\eta\sqrt C))+o(1)\). Choose \(\eta\), then \(C\). This explains the logarithmic-square cutoff and the use of GRH.

\paragraph{Few nearly constant characters.}
Write \(b_\psi=\abs{\sum_{p\in S}\psi(p)}/m\), \(D=\lfloor 1/a\rfloor\le3\). There is a fixed \(\tau>0\) such that the subgroup
\[
 H=\langle\psi:b_\psi>1-\tau\rangle\ \subseteq\widehat G
             \qquad\hbox{has size } d\le D .
\]
We recall the argument of \cite{article,factor}. In the probability inner product on \(\Pp\), the Gram matrix of \(D+1\) distinct characters has operator norm at most \(1+D\epsilon\) by \eqref{eq:pool}. Put \(\alpha=m/N\). If their biases on \(S\) are at least \(w\), the sum of their squared inner products with \(\one_S\) is at least \((D+1)\alpha^2w^2\), and at most \((1+D\epsilon)\alpha\).
Each generator or its inverse is within \(\sqrt{2\tau}\), in normalized \(L^2(S)\), of a unit constant. A word of length at most \(D\) is within \(D\sqrt{2\tau}\) of a unit constant by telescoping, hence has bias at least \(1-D^2\tau\). Since \((D+1)a>1\), choose \(\tau,\epsilon>0\) so that \(1-D^2\tau>0\) and
\[
 (D+1)a(1-D^2\tau)^2>1+D\epsilon .
\]
Fourier characters in the word ball of radius \(D\) therefore number at most \(D\). A word ball gains a new element at every radius until it fills its group, which proves the claim. If \(d=2,3\), let \(\chi\) generate \(H\); it maps \(G\) onto the \(d\)-th roots of unity. If \(d=1\), the quotient below is a point.

\paragraph{Provenance.}
At ProveIt commit \texttt{b6253012cf30}, the vendored \texttt{grh} tree and the upstream tree at \texttt{72fead102f66} both have Git object ID \texttt{2e133df63d7f4da6b3059c4195324f8e9187ba73}. In \texttt{src/Statements.lean}, \texttt{DirichletLZerosOnLine} quantifies over every positive modulus and every Dirichlet character in the open critical strip. In \texttt{src/Main.lean}, \texttt{LFunction\_eq\_zero\_re} proves this statement, while \texttt{LFunction\_ne\_zero\_of\_half\_lt\_re} supplies the half-plane formulation with the principal exception at \(1\) \cite{lean}. The repository attributes the proof to Levent Alp\"oge with Claude. A.~V. Sutherland records a successful independent kernel replay and statement comparison for this revision \cite{lean}. We checked the declarations and Git tree identity and relied on that published record for the build. The downstream proofs are given here on paper.

\section{Lifting from character values to actual residues}
For \(d>1\), let \(\nu\) denote the law of \((\chi(X_1),\ldots,\chi(X_r))\). For \(d=1\), set \(\nu=1\). We include the quotient argument to specify the strength needed to infer \emph{exact} support. Its multivariate coefficient method and lifting principle occur in \cite{factor}.

\begin{lemma}\label{lem:lift}
With \(B\) sufficiently large, uniformly for \(g\in G^r\),
\begin{equation}\label{eq:lift}
 \left|\abs G^{\,r}\pr(X=g)-d^r\nu\bigl(\chi(g_1),\ldots,\chi(g_r)\bigr)\right|\le q^{-2}.
\end{equation}
The support is the complete inverse image of the projected support, with the natural interpretation when \(d=1\).
\end{lemma}
\begin{proof}
Set \(\alpha_i=n_i/m\). For \(\psi_0=1,\psi_1,\ldots,\psi_r\in\widehat G\) the numerator of the partition Fourier coefficient is extracted from
\(\prod_{p\in S}\sum_{i=0}^r z_i\psi_i(p)\). Use \(|z_i|=\alpha_i\) in Cauchy's bound and normalize by the multinomial count. The extra factor is the reciprocal probability of the modal type \((n_0,\dots,n_r)\) in \(m\) independent trials of probabilities \(\alpha_i\), at most \((m+1)^{r+1}\). Indeed changing a type from \(n\) to \(e\) multiplies its probability by \(\prod_i n_i^{e_i-n_i}n_i!/e_i!\le1\).
For phases \(v_i\) of modulus one,
\[
 \left|\sum_i\alpha_i v_i\right|^2
 =1-2\sum_{i<j}\alpha_i\alpha_j(1-\Re(v_i\bar v_j)).
\]
Keeping one pair, applying \(1-x\le e^{-x}\), and multiplying over primes gives, for every \(i\ne j\),
\begin{equation}\label{eq:pair}
 \left|\E\prod_{\ell=1}^r\psi_\ell(X_\ell)\right|
 \le (m+1)^{r+1}\exp\{-m\alpha_i\alpha_j(1-b_{\psi_i\bar\psi_j})\}.
\end{equation}
Choose \(j\) of largest quota. If \((\psi_1,\ldots,\psi_r)\notin H^r\), some \(i\) satisfies \(\psi_i\bar\psi_j\notin H\), because \(\psi_0=1\). Then \(m\alpha_i\alpha_j\ge h/(r+1)\) and the gap is at least \(\tau\). There are at most \(q^r\) tuples of characters. As \(m\le C L^2\), the sum of all these coefficients is at most \(q^{-2}\) for a large enough \(B\). The sum over \(H^r\) in Fourier inversion is exactly \(d^r\nu\) at the projected target.

An occupied projected point has mass at least an inverse polynomial in \(m\) of degree depending only on \(r,d\). Fix an allocation attaining it. Compare any other allocation to this one by drawing an edge from the present to the desired label for each misplaced prime. The directed multigraph is balanced and decomposes into edge-disjoint simple cycles of length at most \(r+1\). Executing a cycle preserves quotas. The projected changes from the cycles multiply to the desired correction in a group of order \(d^r\). In any list of \(d^r\) or more changes two partial products agree, including the empty partial product; remove the intervening block. At most \(d^r-1\) cycles suffice. Thus every allocation is within \(K=(r+1)(d^r-1)\) changed positions of an allocation attaining the target. Such a neighborhood has at most \(\sum_{\ell=0}^K\binom m\ell(r+1)^\ell\) allocations. Its reciprocal exceeds \(q^{-2}\) for large \(q\). This proves that every lift of an occupied class occurs. The projection itself rules out every lift of an empty class.
\end{proof}

\section{The exact allocation problem in two and three classes}\label{sec:allocate}
Fix a character of order two or three, choose \(b\) as above, and call the primes outside its most frequent class exceptional. The reserve makes the individual target conditions particularly simple. The following arguments only require \(n_0\ge 2r\), every recorded quota at least two, and at least \(2r\) primes of class \(b\). These hold under our assumptions for all sufficiently large \(q\).

\paragraph{Order two.}
Each requested minus sign uses at least one exceptional prime. If \(s_\chi<r\) and there are at most \(s_\chi\) requested minus signs, give one exceptional prime to each such block, put the remaining exceptions in the reserve, and fill the quotas with the most frequent class. If \(s_\chi\ge r\), put \(r\) distinct exceptional primes in the reserve, one distinct prime of class \(b\) in every recorded block, and fill the remaining places arbitrarily. Swap the reserved pair for any block whose sign needs to change. The swaps have disjoint supports. This proves the individual and full-coverage tests for the binary projection.

\paragraph{Order three.}
Call an exceptional prime positive or negative according as its value divided by \(b\) is \(\omega\) or \(\omega^{-1}\). Give them costs two and one respectively. A block with target \(\omega\) costs at least two: a single negative prime has the wrong value. A target \(\omega^{-1}\) costs at least one. Reversing the costs gives the other inequality in \eqref{eq:polygon}.

If \(u+v<2r\), all unused exceptional primes will fit in the reserve. When \(u\ge i\) and \(v\ge j\), supply each nontrivial target directly. If \(u<i\), supply \(u\) positive targets directly and use two negative primes for each of the other \(i-u\). The first inequality gives \(v\ge 2(i-u)+j\), which also supplies the negative targets. The case \(v<j\) is symmetric. Fill every unused place with class \(b\).

If \(u+v\ge2r\), put \(2r\) distinct exceptional primes in the reserve and two primes of class \(b\) in each recorded block. Fill the remaining positions arbitrarily. Pair the exceptions in the reserve with the reserved class-\(b\) positions, two pairs per block. For either choice of two nonzero changes in \(\mathbb Z/3\mathbb Z\), their subset sums give the whole group: for equal changes use zero, one or both, and for opposite changes use zero or either singleton. Independent swaps in the blocks give all target vectors. This also covers every case with \(u+v\ge2r\) in \eqref{eq:polygon}. Maximizing \(2i+j\) and \(i+2j\) over \(i+j\le r\) gives \(2r\) in both cases.

\begin{proof}[Conclusion of the proof of Theorem~\ref{th:support}]
Necessity follows by projection under any individual character. If all tests pass, they apply in particular to a generator of \(H\) when \(d>1\). The allocation argument gives that projected target positive mass, and Lemma~\ref{lem:lift} gives the specified tuple of actual units. The same lemma gives every tuple when \(H\) is trivial. For \(a>1/3\), \(d\le2\). Every tuple of character values can be requested, since a character of exact order \(d\) is onto the \(d\)-th roots. This proves the global criteria, including necessity in \eqref{eq:coverage}.
\end{proof}

The binary support has a simple cardinality. For \(a>1/3\) put \(s=\min_{\operatorname{ord}\chi=2}s_\chi\). A quadratic character exists for \(q\ge3\), since the finite abelian unit group has even order. If \(s<r\), a minimizing character has bias \(1-2s/m>1-\tau\) for large \(q\), so generates \(H\). Each sign tuple has \((\abs G/2)^r\) lifts. Including the full-coverage case gives
\begin{equation}\label{eq:hamming}
 \abs{\operatorname{supp}X}=
 \left(\frac{\abs G}{2}\right)^r
       \sum_{j=0}^{\min(r,s)}\binom rj .
\end{equation}
In particular \(s=r-1\) permits every proper marginal to have full support while the joint law misses one entire sign class.

\section{Independent placements and the transition to uniformity}
It remains to quantify the probabilities within the support. For the generator of \(H\) assume \(d=2\) or \(3\), and put \(s=s_\chi\). Write the exceptional values as \(b\xi_1,\ldots,b\xi_s\), with \(\xi_\ell\ne1\). Compare the allocation to the following experiment. Put each of these \(s\) labelled exceptions independently into box \(i\) with probability \(\alpha_i=n_i/m\). Multiply the \(\xi_\ell\)'s in each recorded box, and multiply block \(j\)'s result by \(b^{n_j}\). Denote the resulting law by \(\widetilde\nu\). In the actual partition the exceptions occupy \(s\) distinct slots chosen uniformly without replacement.

\begin{lemma}[Approximation of the projected law]\label{lem:iid}
For fixed \(r\), every Fourier coefficient of \(\nu\) differs from the corresponding coefficient of \(\widetilde\nu\) by \(O_{a,r}(\log(2h)/h)\). Thus in \eqref{eq:lift} one may replace \(\nu\) by \(\widetilde\nu\) at this additive cost on the normalized scale. The exact projected law is still used for deciding support.
\end{lemma}
\begin{proof}
First compare the ordered box labels on \(s\) draws. After \(\ell\) draws without replacement, let \(H_i\) be the occupancy of box \(i\); the next probability is \(p_i=(n_i-H_i)/(m-\ell)\). For any categorical distributions, their relative entropy satisfies
\(\sum p_i\log(p_i/\alpha_i)\le\sum(p_i-\alpha_i)^2/\alpha_i\), by \(\log x\le x-1\). The hypergeometric variance
\(\operatorname{Var}(H_i)=\ell\alpha_i(1-\alpha_i)(m-\ell)/(m-1)\) follows by expanding the indicators of two distinct draws. The chain rule for relative entropy therefore bounds the divergence of the actual label sequence from the independent law by
\[
  \sum_{\ell=0}^{s-1}\sum_{i=0}^r
       \frac{\E(p_i-\alpha_i)^2}{\alpha_i}
       =r\sum_{\ell=0}^{s-1}\frac{\ell}{(m-1)(m-\ell)}
       \ll_r(s/m)^2 .
\]
Here \(s\le2m/3\), since \(b\) is a most frequent class. Pinsker's inequality gives a difference \(O_r(s/m)\) for expectations of functions of modulus at most one. In particular this applies to every Fourier coefficient.

For a second bound, write an exponent tuple \((k_0,\ldots,k_r)\in(\mathbb Z/d\mathbb Z)^{r+1}\) with \(k_0=0\). In a nontrivial coefficient some exponent differs from the exponent of a block of largest quota. In \eqref{eq:pair} the corresponding ratio is \(\chi\) or its inverse. For \(d=2\) its gap is \(2s/m\). For \(d=3\),
\[
 1-b_\chi^2=3\{(m-s)s+uv\}/m^2\ \ge\ s/m,
 \qquad 1-b_\chi\ge s/(2m).
\]
It follows that the actual nontrivial coefficient has modulus at most
\((m+1)^{r+1}\exp(-hs/(2(r+1)m))\).

For the independent experiment put \(\zeta=e^{2\pi i/d}\) and \(z=\sum_{i=0}^r\alpha_i\zeta^{k_i}\). Each exceptional prime contributes \(z\) or \(\bar z\), up to the common deterministic block factors. Its coefficient therefore has modulus \(\abs z^s\). For \(d=2\) the largest possible \(\abs z\) is \(1-2t\). Indeed both sign classes have total weight at least \(t\), and isolating a smallest block gives equality. For \(d=3\) write \(P_0,P_1,P_2\) for the total weights of the three phases. Then
\[
 |z|^2=1-3(P_0P_1+P_1P_2+P_2P_0)\le1-3t(1-t).
\]
To see the inequality, one nonempty phase class has mass \(p\in[t,1/2]\), and the sum of pair products is at least \(p(1-p)\). Again isolate a smallest block for equality. If the smallest block is the reserve, assign an equal nonzero exponent to every other block. Thus the independent coefficient is at most \(\exp(-c ts)\) for an absolute \(c>0\).

Now \(m\le C L^2\ll_{a,r} h^2\). When \(s/m\le K_r\log(2h)/h\), the entropy bound suffices. Otherwise, on choosing \(K_r\) large, both Fourier coefficients are \(O_{a,r}(h^{-1})\) by the exponential estimates. Fourier inversion on the group of size \(d^r\) proves the assertion about normalized point masses.
\end{proof}

\begin{proof}[Proof of Theorem~\ref{th:mix}]
For any character outside \(H\) of order two or three, comparison with its most frequent constant gives \(b_\psi\ge1-2s_\psi/m\). Thus \(s_\psi\ge\tau m/2\), and \(B_{\operatorname{ord}\psi}(t)^{s_\psi}=O(e^{-c\tau h})\). Inside \(H\), the maximum modulus of a nontrivial independent coefficient is precisely \(B_d(t)^s\), attained as in Lemma~\ref{lem:iid}. Conjugate cubic generators have the same \(s\). The maximum nontrivial coefficient of the actual projected law therefore agrees with \(\mathcal M_q\) to the stated additive error, with a zero maximum for a trivial quotient.

For a nontrivial character its Fourier coefficient is the uniform average of \((\abs{G}^{\,r}\pr(X=g)-1)\) times a unit phase. Its modulus is at most \(D_q\). Conversely, Fourier inversion bounds \(D_q\) by the sum of the \(d^r-1\) nontrivial in-quotient moduli and the \(q^{-2}\) bound for all other characters. Use \(d\le3\), or \(d\le2\) for \(a>1/3\). In the one-product quadratic case there is at most one nonconstant in-quotient coefficient, proving the identity.
\end{proof}

\paragraph{One product at finer precision.}
For \(a>1/3\), let \(k=n_1\) and choose a quadratic character minimizing \(s=s_\chi\). The exact parity coefficient for the number of exceptions in the product is
\[
 P_{m,s}(k)=\binom{m}{s}^{-1}\sum_{j=0}^s(-1)^j
                 \binom{k}{j}\binom{m-k}{s-j},
\]
with out-of-range binomial coefficients zero. At \(y=C(\log q)^2\),
\(D_q=|P_{m,s}(k)|+O(q^{-2})\). In fact if \(H\) has order two its generator uniquely maximizes the quadratic bias, and its contribution changes sign between the two fibres. If \(H\) is trivial every nontrivial coefficient is bounded by the total \(q^{-2}\) tail. This gives an integer formula even in the balanced regime where the independent approximation can vanish exactly.

The same one-product mixing criterion holds for \(y=L^{F}\), \(F>2\) fixed, throughout the shorter range
\[
  h\ge c\,L/\log L,\qquad c>2(F-1)/(F(F-2)).
\]
Here \texttt{dir:thm:short} of \cite{article} bounds the out-of-quotient Fourier sum by \(\exp(-\gamma L/\log L)\) for some \(\gamma>0\). Its bias lemma gives the same order-two bound for the exceptional subgroup when \(a>1/3\). The exact parity formula therefore has this new remainder in place of \(O(q^{-2})\). In Lemma~\ref{lem:iid}, \(m\le L^F\) remains polynomial in \(h\); its proof gives the same \(O_{a,F,c}(\log(2h)/h)\) approximation. This addresses both prime cutoffs appearing in the article's half-density question.

\paragraph{A continuum of limiting biases at half retention.}
Take odd prime \(q\), fix \(a\in(1/3,1/2]\) and \(m=\lceil N/2\rceil\). Choose a most frequent quadratic sign in \(\Pp\). At least \(m\) primes have that sign, and \eqref{eq:pool} supplies a positive proportion of the other sign. Take \(k=\lceil B\log q\rceil\) recorded places, with \(h=k\) for large \(q\). Given any fixed \(\lambda\ge0\), choose exactly \(s=\lfloor\lambda m/k\rfloor\) opposite-sign primes and \(m-s\) of the first sign. Here \(k\to\infty,\ k/m\to0\), and \(sk/m\to\lambda\). Theorem~\ref{th:mix} gives the sharp limit
\[
                      D_q\longrightarrow e^{-2\lambda}.
\]
For \(\lambda>0\) all classes are eventually reached and the two normalized class counts tend to \(1+e^{-2\lambda}\) and \(1-e^{-2\lambda}\), up to their sign labelling. This shows why any criterion based only on retaining a fraction tending to \(1/2\) misses a full range of limiting laws.

\paragraph{Sharpness at one-quarter retention.}
Take arbitrarily large primes \(q\equiv5\pmod{12}\). There is one quadratic character, no character of order three, and a character \(\theta\) of order four. Retain a most frequent quartic class \(c\) and exactly \(r\) primes from the class \(ic\). The first class contains at least \(N/4\) primes. The second contains \(\gg N\), by expanding its indicator using \(1,\theta,\theta^2,\theta^3\) in \eqref{eq:pool} with \(\epsilon<1/4\). The unique quadratic character has \(r\) minority primes, so every test in \eqref{eq:coverage} passes. Prescribe \(\theta(g_j)c^{-n_j}=-1\) in every recorded block. This needs at least two exceptional primes per block, exceeding the supply. All quotas can be taken of order \(m/(r+1)\). Thus the quadratic and cubic coverage tests cease to be sufficient at \(a=1/4\).

\section{How many products can be prescribed at once?}\label{sec:growing}
The article asks how fast the number of factors can grow at retained density at most one half. For one common modulus and an unrestricted residue in the reserve block, the answer has the same order as the elementary capacity bound. The result in this section allows every fixed density \(0<a\le1\). All constants are independent of the number of target blocks.

\begin{theorem}[Growing number of targets at any positive density]\label{th:growing}
Fix \(0<a\le1\) and \(A>0\). There are \(C=C(a)\), \(B=B(a,A)\) such that the following holds under Dirichlet GRH for every sufficiently large \(q\). Take any \(S\subseteq\Pp\) of size \(m\ge aN\) and any \(r\ge1\) for which the quotas, including the reserve, satisfy \(h\ge B\log q\). Put \(D=\lfloor1/a\rfloor\). There is a group \(H\subseteq\widehat G\) of size \(d\le D\), depending only on \(q,S\), with evaluation quotient \(\pi:G\to\widehat H\). Write \(\nu\) for the joint law of the \(r\) projected products. The support is exactly the inverse image of \(\operatorname{supp}\nu\), and at each occupied target,
\[
  \pr(X=g)=\frac{d^r}{|G|^r}\nu(\pi^r(g))(1+\theta_g),
       \qquad |\theta_g|\le q^{-A}.
\]
For \(a>1/4\), the target tests of Theorem~\ref{th:support} and the coverage thresholds \eqref{eq:coverage} hold throughout this growing range. Balanced quotas allow \(r\asymp_{a,A} \log q/\log\log q\) even at the least allowed retained density. The number of primes available makes this order maximal under the stipulated quota size.
\end{theorem}

\begin{proof}
The Gram argument after \eqref{eq:pool} works verbatim with any fixed \(D\). Choose \(\tau=\tau(a)\), then \(C\), so the subgroup generated by biases exceeding \(1-\tau\) has size \(d\le D\). Put \(k=r+1\). For \(B\ge1\), \(k\le m/\log q\). From \eqref{eq:pool} and \(m\le N\) there is a fixed \(K=3(C+1)\) such that
\begin{equation}\label{eq:growsize}
 k\log(m+1)\le K\log q .
\end{equation}
Indeed \(N\sim C(\log q)^2/\log(C(\log q)^2)\). We give two estimates whose losses remain bounded by fixed powers of \(q\).

\emph{The Fourier error.}
In the Cauchy estimate leading to \eqref{eq:lift}, retain every pair of labels whose characters lie in different cosets of \(H\). Give each coset the weight \(w_c=\sum_{i:\psi_i\in c}n_i/m\), including \(\psi_0=1\). With \(c_*\) a coset of largest weight and \(t_*\) the number of labels outside it,
\[
 \sum_{c<c'}w_cw_{c'}=\tfrac12(1-\textstyle\sum_c w_c^2)
 \ge \tfrac12(1-w_{c_*})\ge t_* h/(2m).
\]
Thus the Fourier coefficient is bounded by \((m+1)^k e^{-\tau t_* h/2}\). A tuple outside \(H^r\) has \(t_*\ge1\), since the reserve character is trivial. There are at most \(q\binom{k}{t_*}d^k q^{t_*}\) tuples for a given \(t_*\): choose the dominant coset, the other labels, and their characters. This count permits every possible reserve character and hence majorizes the required count. If \(B\ge8/\tau\), then \(kq e^{-\tau h/2}\le q^{-2}\) for large \(q\). Summing the binomial bound, using \((1+x)^k-1\le2kx\) for \(kx\le1/2\), bounds the total error by
\begin{equation}\label{eq:growerror}
 (m+1)^k q d^k\big[(1+q e^{-\tau h/2})^k-1\big]
 \le 2q^{2K+3-\tau B/2}.
\end{equation}
Here \(k\le q\) and \(d\le m+1\) for large \(q\). Orthogonality on \(H^r\) gives the principal term \(d^r\nu(\pi^r(g))\) in the normalization of \eqref{eq:lift}.

\emph{Repair of the color counts.}
Call \(\pi(p)\in\widehat H\) the color of a prime, and let \(e\le d\) be the exponent of \(\widehat H\). Fix one allocation witnessing an occupied projected tuple. Its \(k\)-by-\(d\) table \(Y\) records the number of each color in each block. The table \(X'\) of any other allocation has the same row and column sums. We construct a table \(X''\) with those margins and with every entry congruent to the corresponding entry of \(Y\) modulo \(e\). This ensures that each block has exactly the witnessed color product.

Let \(\sigma_{ij}\in\{0,\ldots,e-1\}\) be the residue of \(Y_{ij}\). Take the row and column sums of \((Y-\sigma)/e\) as target margins. Start at the nonnegative integer table
\[
 Z_{ij}=\max\{0,\lfloor(X'_{ij}-\sigma_{ij})/e\rfloor\}.
\]
Initially \(\|\sigma+eZ-X'\|_1\le e k d\), where the norm sums absolute entries. Both the total absolute row discrepancy and the total absolute column discrepancy of \(Z\) from the target margins are at most \(kd\), by the equal margins of \(X',Y\). Delete units from overfull rows until none remains, and then from overfull columns until none remains. Each stage costs at most \(kd\), as row deletions can only decrease column excesses. The remaining summed row deficit equals the summed column deficit and is at most \(3kd\). Fill it by adding units at intersections of deficient rows and columns. Writing \(Z_*\) for the repaired table, we obtain
\[
 X''=\sigma+e Z_*,\qquad \|X''-X'\|_1\le 6e k d\le6D^2 k.
\]
For each color, transfer individual primes from its surplus blocks to its deficient blocks. This changes at most \(6D^2 k\) prime labels, preserves every quota in the final allocation, and reaches the witnessed projection. The number of allocations within that many changes of any one allocation is at most \((1+mk)^{6D^2k}\): pad a list of edits with empty moves. Since all allocations have the same probability, the neighborhoods of the allocations realizing the projection cover the full space. Thus
\[
  \nu(\pi^r(g))\ge (1+mk)^{-6D^2 k}\ge q^{-12D^2 K}
\]
at every occupied point by \eqref{eq:growsize}. Take, for instance,
\(B\ge (2/\tau)(A+12D^2K+2K+6)\), also with \(B\ge1\). Dividing \eqref{eq:growerror} by the occupied main term gives the stated relative error. An empty projection rules out every lift by definition.

Finally, \eqref{eq:pool} and \(kh\le m\) imply \(r=O_{a,A}(\log q/\log\log q)\), hence \(h\ge2r\) uniformly for large \(q\). When \(a>1/4\) the quotient has order at most three. The finite sign and cubic allocation arguments of Theorem~\ref{th:support} apply unchanged: the quotas can hold \(2r\) exceptional primes and a most frequent class supplies at least \(m/3\ge2r\) anchors. Balanced quotas with \(k=\lfloor aN/(2B\log q)\rfloor\) demonstrate the claimed attainable order for every permitted retained set.
\end{proof}

The quotient in this theorem uses a single unit group \(G\). The next section allows an independently chosen modulus for each block.

\section{A common prime list for growing families of different moduli}\label{sec:mixed}
We now allow a separate modulus for every block and use every retained prime. A block of modulus one plays the role of a reserve. The open question on growing families following \texttt{dir:thm:minority} asks whether relative counts on occupied fibres survive below majority retention. They survive up to the full order of growth allowed by the prime cutoff and the quotas.

\begin{theorem}[Independent moduli after positive-density retention]\label{th:mixed}
Fix \(0<a\le1\) and \(A>0\), and put \(D=\lfloor1/a\rfloor\). There are positive constants \(B,K_1,K_2,c\), depending only on \(a,A\), with the following property under Dirichlet GRH. Write \(L=\log Q\). For every sufficiently large \(Q\) and every integer \(2\le R\le cL/\log L\), there is a single list \(\mathcal S\) of between \(K_1 R L\) and \(2K_1 R L+2\) distinct primes in \([x,2x]\), where \(x\asymp_a L^2\) and \(x\le K_2 L^2\).

Choose \(2\le r\le R\) and \(1\le q_j\le Q\) independently. Set \(M=\operatorname{lcm}_j q_j\), \(g_0=\gcd_j q_j\), and \(U=\{p\in\mathcal S:p\nmid M\}\). Retain any \(T\subseteq U\) of size \(m\ge a|U|\). For a uniform partition of \(T\) into \(r\) labelled blocks of sizes \(n_j\) with \(h=\min_j n_j\ge BL\), let \(\mu\) be the law of their product vector in \(V=\prod_j(\mathbb Z/q_j\mathbb Z)^\times\). Balanced admissible quotas exist for each such \(T\).

There is a subgroup \(E\subseteq\widehat V\), independent of the quotas, containing the diagonal character group \(\Delta\simeq\widehat{(\mathbb Z/g_0\mathbb Z)^\times}\), and
\[
 [E:\Delta]\le D^{r-1}.
\]
For evaluation \(\Pi:V\to\widehat E\), let \(\nu=\Pi_*\mu\). The support of \(\mu\) is precisely \(\Pi^{-1}(\operatorname{supp}\nu)\). On it,
\[
   \mu(v)=\frac{|E|}{|V|}\nu(\Pi(v))(1+\theta_v),
       \qquad |\theta_v|\le Q^{-A}.
\]
The projected support has at most \([E:\Delta]\) points. For \(a>1/2\) the sole constraint is that the product of the target residues agree with \(\prod_{p\in T}p\) modulo \(g_0\).
\end{theorem}

\begin{proof}
\emph{A list that can test short products of characters.}
Choose \(\tau,\eta>0\) so that \(1-D^2\tau>0\) and
\((D+1)a(1-D^2\tau)^2>1+D\eta\), and set \(\delta=\eta/4\), reducing \(\eta\) if needed to be less than one. We use the source's common-list version of the prime-pool estimate \eqref{eq:pool} (the one-set lemma in \cite{factor}, labelled \texttt{dir:lem:oneset} in the enlarged article). With parameter \(Q_*=Q^{2D}\) it supplies, for every integer length parameter
\[
  C_\delta\log(2Q_*)\le N_*\le c_\delta\log^2(2Q_*)/\log\log(3Q_*),
\]
a single list of between \(N_*\) and \(2N_*\) primes in \([x,2x]\), \(x=K_\delta\log^2(2Q_*)\), on which every nonprincipal character of modulus at most \(Q_*^2=Q^{4D}\) has mean of modulus at most \(\delta\). The source derives the dyadic estimate by the smoothed GRH explicit formula and selects the list by independent sampling.

Take \(B\ge \max\{1,(2/\tau)(A+36D^2+13)\}\). Next choose \(K_1\) large enough for the lower bound on \(N_*=\lceil K_1RL\rceil\) and \(aK_1/2\ge 2(B+1)\). Finally choose \(0<c\le1\) small enough for the upper bound throughout the range of \(R\). All choices are independent of \(q_j,T,r\). At most \(rL/\log x\) list primes divide \(M\), as follows by taking their product. This is \(o(|\mathcal S|)\) uniformly, so every indicated character has mean on \(U\) of modulus at most \(\eta\) for large \(Q\), and \(m\ge a K_1 R L/2\). This proves the balanced-quota assertion. Also, since \(r\le L/\log L\) and \(m\le 2K_1 R L+2\), for large \(Q\) one has
\begin{equation}\label{eq:mixedsize}
            r\log(m+1)\le 3L .
\end{equation}
The coefficient three here is independent of the fixed choice of \(K_1\).

\emph{The subgroup seen by the available characters.}
Inflate characters of each block to the unit group modulo \(M\), writing the inflated character as \(\widetilde\psi_j\). Generate \(J\) by all ratios \(\widetilde\psi_i\overline{\widetilde\psi_j}\) of block characters whose bias on \(T\) exceeds \(1-\tau\), allowing all choices of the two characters. To control this subgroup, take two different words of length at most \(D\) in the generators and inverses, viewed as different characters. Their quotient involves at most \(4D\) block characters counting repetitions. It is induced by a nonprincipal character modulo the least common multiple of the participating moduli, an integer at most \(Q^{4D}\). On \(U\) the induced and inflated values coincide. Thus the corresponding distinct word characters have Gram off-diagonal entries bounded by \(\eta\), even though \(M\) can be larger than any fixed power of \(Q\). Each word has bias at least \(1-D^2\tau\) on \(T\) by the telescoping estimate used after \eqref{eq:pool}. The same Gram-matrix and word-ball argument with \(|T|/|U|\ge a\) now gives \(|J|=d\le D\).

Define \(E\) to consist of tuples \((\psi_1,\dots,\psi_r)\) with all inflated characters in a single coset of \(J\). This is a subgroup. Taking their ratios with \(\widetilde\psi_1\) gives a homomorphism \(E\to J^{r-1}\) whose kernel consists of equal inflated characters. A common character descends to the gcd of the moduli: its unique primitive conductor divides every \(q_j\). Conversely every character modulo \(g_0\) inflates to such a tuple. This identifies the kernel with \(\Delta\) and proves the index bound, including unit moduli.

\emph{Uniform Fourier decay.}
For a character tuple outside \(E\), let \(t\ge1\) be the number of labels outside a coset of largest combined quota. Distinct cosets have bias gap at least \(\tau\) on all ratios of their corresponding labels, by construction of \(J\). The pair estimate retaining all cross-coset pairs in Theorem~\ref{th:growing} bounds that coefficient by \((m+1)^r e^{-\tau t h/2}\). There are at most
\[
 rQ\binom rt d^r Q^t
\]
such tuples for each \(t\): specify one character in the dominant coset from some block, the outside labels, and their characters. Each inside label has at most \(d\) possibilities since inflation from its modulus is injective. Since \(rQ e^{-\tau h/2}\le1/2\) for large \(Q\), summing as in \eqref{eq:growerror} gives
\[
 \sum_{\boldsymbol\psi\notin E}|\widehat\mu(\boldsymbol\psi)|
 \le (m+1)^r rQ d^r\big[(1+Qe^{-\tau h/2})^r-1\big]
 \le 2Q^{10-\tau B/2}.
\]
We used \eqref{eq:mixedsize}, \(d\le m+1\), and \(r\le Q\). Fourier inversion over \(E\) contributes exactly \(|E|\nu(\Pi(v))\) to \(|V|\mu(v)\).

\emph{Mass of an occupied projection.}
Color each prime by its value in \(\widehat J\) under evaluation. For a tuple in \(E\), put \(\delta_j=\widetilde\psi_j/\widetilde\psi_1\in J\). If the blocks have integer products \(P_j\), the value of the tuple is
\[
 \widetilde\psi_1\Big(\prod_{p\in T}p\Big)\prod_{j=1}^r\delta_j(P_j).
\]
Consequently preserving the \(r\) color products of any witness preserves its complete \(E\)-projection. The integer-table repair in Theorem~\ref{th:growing}, with \(r\) rows and \(d\) colors, realizes the witness's color products from every allocation using at most \(6D^2r\) edited prime labels. Hence each occupied projection has mass at least
\[
    (1+mr)^{-6D^2r}\ge Q^{-36D^2}
\]
by \eqref{eq:mixedsize}. The Fourier error is at most \(Q^{-A}\) times every occupied main term for large \(Q\) by our choice of \(B\). Empty projections cannot occur. All occupied projections restrict to the same value on \(\Delta\), determined by the total product. This gives the bound on their number and the assertion for \(D=1\).
\end{proof}

The order of the growing range is sharp for the stated size of primes and balanced-quota guarantee. Even with no deletions, \(r\) blocks of size \(B\log Q\) use at least \(rB\log Q\) distinct primes, whereas \([1,2K_2(\log Q)^2]\) contains only \(O_{K_2}((\log Q)^2/\log\log Q)\) primes. Thus the loss allowed by the quotient introduces no further restriction in order of magnitude on the number of independently prescribed moduli.


\section{Finite checks and scope}
The supplementary Python program counts exact residue tuples by processing each prime in turn. The choices are the reserve and the two recorded blocks. Its state records both used quotas and both residues. A second computation enumerates allocations directly on a smaller example. The following finite instances illustrate the support transitions.

\begin{center}
\begin{tabular}{@{}llllr@{}}
\toprule
Modulus & Class data & Quotas (reserve first) & Occupied classes & Residue pairs\\
\midrule
\(11\) & \(s=0\) & \(6,6,6\) & \(1/4\) & \(25/100\)\\
       & \(s=1\) & \(6,6,6\) & \(3/4\) & \(75/100\)\\
       & \(s=2\) & \(6,6,6\) & \(4/4\) & \(100/100\)\\[2pt]
\(13\) & \((u,v)=(0,0)\) & \(5,5,5\) & \(1/9\) & \(16/144\)\\
       & \((u,v)=(1,1)\) & \(5,5,5\) & \(7/9\) & \(112/144\)\\
       & \((u,v)=(2,0)\) & \(5,5,5\) & \(6/9\) & \(96/144\)\\
       & \((u,v)=(2,1)\) & \(5,5,5\) & \(9/9\) & \(144/144\)\\
\bottomrule
\end{tabular}
\end{center}

Modulo \(11\), keep the smallest \(18-s\) residue primes and \(s\) nonresidue primes up to \(199\), excluding \(11\). For \(s=2\) there are \(17\,153\,136\) allocations. The \(++\) sign class gets \(5\,045\,040\), while each of the other three gets \(4\,036\,032\). Modulo \(13\) use the smallest requested numbers of primes up to \(239\) in the three classes \(p^4=1,3,9\bmod 13\), with counts \(15-u-v,u,v\). Two exceptions of opposite types give strictly larger support than two of the same type.

The independent finite urn checks cover \(2\,464\) binary cases, \(15\,072\) binary Fourier coefficients and \(104\,505\) cubic target vectors from \(4\,231\) capacity and inventory profiles. They check the combinatorial formulas; the proof establishes the uniform arithmetic conclusions.

A separate implementation of the table repair checks \(28\,183\) ordered pairs of tables with equal margins, together with larger deterministic examples. It verifies the required congruences, margins and edit bounds at each stage. A direct enumeration with moduli \(15,21\) and quotas \(7,7\) reaches precisely the \(48\) out of \(96\) unit pairs obeying the common congruence \(g_1g_2\equiv2\bmod3\). The \(3\,432\) assignments give \(65\) to \(80\) representations per occupied pair.

\paragraph{Related work and priority.}
Booker and Pomerance studied squarefree smooth representatives of individual residue classes \cite{BP}. Klurman, Shparlinski and Ter\"av\"ainen and subsequently Hughes obtained common short prime lists for Boolean subset products for almost all moduli \cite{KST,Hughes}. Pach treats the distribution of fixed-cardinality subset sums of an entire abelian group \cite{Pach}. The closest sources for our setup are the repository's deletion theorem and its joint-factorization theorem \cite{article,factor,merged}. The latter already contains a finite-quotient lifting result and a three-block example with a single nonresidue. The target inequalities, the mixing criterion, and the growing-factor argument at arbitrary positive density proved here extend those results. Historical priority beyond the sources examined is undetermined.

\paragraph{Further questions.}
At density \(1/4\) the group of order four can be cyclic or have exponent two. What are the exact support inequalities in those two cases? For fixed density, can the logarithmic quota for a prescribed relative error be given with its optimal leading constant? For growing \(r\), which statistics determine the approach to uniformity \emph{across} the occupied quotient classes? The growing-factor theorem controls the relative counts within each class.

\begin{thebibliography}{10}
\setlength{\itemsep}{1pt}\setlength{\parsep}{0pt}
\footnotesize\raggedright
\bibitem{lean} \texttt{lsernau/grh-lean}, \emph{GRH for Dirichlet \(L\)-functions, in Lean 4}, readable edition, upstream commit \texttt{72fead102f66}. \repo{vendor/grh-lean/README.md}{Vendoring record},
\repo{vendor/grh-lean/grh/src/Statements.lean}{statements} and
\repo{vendor/grh-lean/grh/src/Main.lean}{theorem exports}. A.~V. Sutherland, \href{https://math.mit.edu/~drew/shortgrh/}{verification files for the shorter Lean proof}.
\bibitem{article} V.~Reshetnikov, ProveIt collection, \emph{Beyond the Riemann hypothesis for Dirichlet \(L\)-functions}, revision \texttt{b6253012cf30}. \repo{Article/README.md}{Article overview} and \repo{Article/sections/dir.tex}{direct arithmetic consequences}. See \texttt{dir:lem:bias}, \texttt{dir:thm:quotient} and the question on density near one half.
\bibitem{factor} ProveIt research report, \emph{Simultaneous congruences for factors of a squarefree integer}, revision \texttt{b6253012cf30}. \repo{Research/squarefree_factorizations/factorizations.tex}{Full source}. See the results on prescribed partitions and positive retained density.
\bibitem{merged} ProveIt collection, expanded 120-report revision \texttt{3d8689ceae76}. \href{https://github.com/VladimirReshetnikov/ProveIt/blob/3d8689ceae76501a5c50f67da37e85369aac4fb6/Analysis/GeneralizedRiemannHypothesis/Article/sections/dir.tex}{Merged factorization theorem and open questions}. See \texttt{dir:thm:minority} and the closing questions on half density and growing numbers of factors.
\bibitem{MV} H.~L. Montgomery and R.~C. Vaughan, \emph{Multiplicative Number Theory I. Classical Theory}, Cambridge Studies in Advanced Mathematics 97, Cambridge University Press, 2007, Ch.~13.
\bibitem{BP} A.~R. Booker and C.~Pomerance, \emph{Squarefree smooth numbers and Euclidean prime generators}, Proc. Amer. Math. Soc. 145 (2017), 5035--5042. \href{https://arxiv.org/abs/1607.01557}{arXiv:1607.01557}.
\bibitem{KST} O.~Klurman, I.~E. Shparlinski and J.~Ter\"av\"ainen, \emph{On the generation of multiplicative groups by small primes}, 2026. \href{https://arxiv.org/abs/2609.25393}{arXiv:2609.25393}.
\bibitem{Hughes} S.~D. Hughes, \emph{Few small primes suffice to cover unit groups}, 2026. \href{https://arxiv.org/abs/2610.04678}{arXiv:2610.04678}.
\bibitem{Pach} P.~P. Pach, \emph{Equidistribution of subset sums}, 2025. \href{https://arxiv.org/abs/2505.12496}{arXiv:2505.12496}.
\end{thebibliography}
\end{document}
